\section{Conclusion}
\label{sec:conclusion}

This paper proposed an AST-native code surface for multi-agent LLM
workflows comprising a unified Markdown outline emitter (with three
discovery surfaces), a discipline subsystem that makes outline-first
reading the default for files above a configurable LoC threshold,
and a plan/execute pair contract for symbolic surgery with a
compile-gated execute step satisfying a false-negative-only
correctness contract.

The measured cross-language compression ratios (Go 82.0\%, Python
84.1\%, JavaScript 83.0\%, TypeScript 75.2\%) on a realistic
benchmark fixture establish the read-side mechanism's floor. The
end-to-end token reduction on real agent traces, the adoption
trajectory of the discipline subsystem, and the false-negative rate
of the compile gate are the pre-registered claims whose datasets
are still accumulating; we are transparent about which numbers are
measured and which are owed.

The architectural commitments --- one emitter, three surfaces; the
hook as the discipline mechanism, not a discoverable feature; the
gate as a hard precondition for symbolic surgery, not an optional
flag --- are the contribution we stake the design on. They compose
with the model-tier-selection advisor from companion work
\cite{blackrim_moa_paper} multiplicatively: outline-first reading
shrinks the prompt; tier-selection picks the cheapest model that
can serve the (smaller) prompt. The savings stack.

The methodological stance is deliberate: data first, design
second. If the discipline subsystem fails to drive adoption above
the 80\% threshold, it does not get promoted to block-mode. If the
compile gate's false-negative rate exceeds the conservatism budget,
the threshold relaxes. If the pattern-DSL fidelity turns out to be
unacceptable for production intents, the YAML-authoring path stays
optional and the in-tree Go matcher remains the load-bearing
backend. The architecture is committed; the parameters bend to the
data.
